%! Modeling with Exponential Functions: Interest, Half-Life, and e — Unit 6, Lesson 6.4 — Homework (Answer Key)
\documentclass[10pt]{article}
\usepackage{algebra2-article}
\usepackage{algebra2-key}   % pulls in -boxes; do NOT also load -boxes
\newcommand{\TallMath}[1]{$\displaystyle #1\rule[-1.4em]{0pt}{3.2em}$}
% ln(1.06) = 0.0582689

\begin{document}

\pageheader{Unit 6, Lesson 6.4}{Homework --- Answer Key}
\namedateperiod

\vspace{0.10in}

\begin{notesbox}{Practice}
Every answer needs a visible substitution. Round money to the nearest cent and milligrams to two
decimal places.

\begin{enumerate}[leftmargin=1.9em, itemsep=11pt]
  \item \textbf{Four accounts.} Write the base and the exponent first, then evaluate.
    \begin{center}
    \small
    \renewcommand{\arraystretch}{1.8}
    \begin{tabularx}{\linewidth}{@{}l X c c c@{}}
    \hline
    & \textbf{Deposit, rate, compounding, time} & \textbf{base} & \textbf{exponent} & \textbf{$A$} \\
    \hline
    \textbf{(a)} & \$1200 at $4.5\%$, quarterly, $6$ years & \ans{$1.01125$} & \ans{$24$} & \ans{\$1569.59} \\
    \textbf{(b)} & \$7500 at $2.4\%$, monthly, $3$ years & \ans{$1.002$} & \ans{$36$} & \ans{\$8059.34} \\
    \textbf{(c)} & \$2500 at $5\%$, annually, $12$ years & \ans{$1.05$} & \ans{$12$} & \ans{\$4489.64} \\
    \textbf{(d)} & \$900 at $6\%$, daily ($n=365$), $4$ years & \ans{$1+\frac{0.06}{365}$} & \ans{$1460$} & \ans{\$1144.10} \\
    \hline
    \end{tabularx}
    \end{center}
    \emph{Work space --- show the substitution for at least two of the four.}
    \vspace{1.05in}

  \item \textbf{Quarterly against continuous.} \$6000 is invested at $4\%$ for $10$ years.

    \vspace{3pt}
    \hspace*{1.0em} Compounded quarterly: $A=$ \ans{$6000(1.01)^{40}$} $=$ \ans{\$8933.18}

    \vspace{3pt}
    \hspace*{1.0em} Compounded continuously: $A=$ \ans{$6000e^{\,0.4}$} $=$ \ans{\$8950.95}

    \vspace{3pt}
    Over ten years, continuous compounding is worth \ans{\$17.77} more. Explain why that is so small
    when the compounding went from $40$ times to infinitely many. \ansline{Each extra chop of the year adds less than the one before, and the balances close in on the ceiling $Pe^{rt}$ --- ``infinitely often'' is barely better than quarterly.}

  \item \textbf{Caffeine.} A $200$ mg coffee has a half-life of about $5$ hours in the body.

    \vspace{3pt}
    \hspace*{1.0em} $A(t)=$ \ans{$200\left(\frac12\right)^{t/5}$} \; where $t$ is in \ans{hours}.

    \vspace{3pt}
    \begin{center}
    \small
    \renewcommand{\arraystretch}{1.4}
    \begin{tabular}{|l|c|c|c|c|c|}
    \hline
    $t$ (hours) & $0$ & $5$ & $10$ & $15$ & $20$ \\ \hline
    $A(t)$ (mg) & $200$ & \ans{100} & \ans{50} & \ans{25} & \ans{12.5} \\ \hline
    \end{tabular}
    \end{center}

    \vspace{2pt}
    \textbf{(a)} How much is left $8$ hours after the coffee? \ans{65.98 mg} \; Check it against the
    table: it must land between \ans{100 mg} and \ans{50 mg}.

    \vspace{3pt}
    \textbf{(b)} According to this model, is the caffeine ever completely gone? \ans{No} \;
    Name the graph feature that decides it. \ans{the horizontal asymptote $A=0$; the range is $(0,200]$}

  \item \textbf{Read the doubling off the graph.} \$800 is invested at $6\%$ compounded annually, so
    $A(t)=800(1.06)^{t}$.
    \begin{center}
    \begin{minipage}[c]{0.44\linewidth}
    \centering
    \begin{tikzpicture}[x=0.25cm, y=0.0019cm]
      \draw[linegray!40, very thin, xstep=2, ystep=400] (0,0) grid (20,2400);
      \draw[->, semithick, charcoal] (0,0)--(21.5,0) node[above right, font=\scriptsize]{$t$ (yr)};
      \draw[->, semithick, charcoal] (0,0)--(0,2600) node[left, font=\scriptsize]{$A$};
      \foreach \t in {5,10,15,20} \node[below, font=\tiny] at (\t,0) {$\t$};
      \foreach \a in {800,1600,2400} \node[left, font=\tiny] at (0,\a) {$\a$};
      \draw[very thick, forest, domain=0:20, samples=120, smooth]
        plot (\x,{800*exp(0.0582689*(\x))});
      \draw[goldacc, very thick] (0,1600)--(20,1600);
      \node[goldacc, font=\scriptsize, above] at (4,1620) {$A=1600$};
    \end{tikzpicture}
    \end{minipage}
    \begin{minipage}[c]{0.52\linewidth}
    \small
    \textbf{(a)} $A(11)=$ \ans{\$1518.64} \quad $A(12)=$ \ans{\$1609.76}

    \vspace{3pt}
    \textbf{(b)} The money doubles between year \ans{11} and year \ans{12}.

    \vspace{3pt}
    \textbf{(c)} Write the equation you would have to solve to find the exact doubling time, then
    divide to simplify it. \ans{$800(1.06)^{t}=1600$, so $(1.06)^{t}=2$}

    \vspace{3pt}
    \textbf{(d)} Try Step 1 of Lesson 6.3 on it. What stops you? \ans{$2$ is not a power of $1.06$ --- there is no common base}
    \end{minipage}
    \end{center}

  \item \textbf{Error analysis.} Name the error, then give the correct model and value.

    \vspace{2pt}
    \textbf{(a)} Tobias models \$5000 at $6\%$ compounded monthly for $4$ years as
    $A=5000(1.06)^{48}$, and reports \$81{,}969.36.
    \ansline{He counted the periods right ($48$ months) but never split the rate: his base pays a full year's interest every month.}
    \ansline{Correct: $A=5000(1.005)^{48}=\$6352.45$. The size of his answer is the tell --- \$5000 does not become \$82{,}000 in four years.}

    \vspace{2pt}
    \textbf{(b)} Marisol models a $90$ mg dose with a $6$-hour half-life as
    $A(t)=90\left(\frac12\right)^{t}$, and reports that $0.00034$ mg is left after $18$ hours.
    \ansline{Her exponent counts \emph{hours}, not half-lives, so her model halves the dose $18$ times instead of $3$.}
    \ansline{Correct: $A(t)=90\left(\frac12\right)^{t/6}$, so $A(18)=90\left(\frac12\right)^{3}=11.25$ mg.}
\end{enumerate}
\end{notesbox}

\begin{extensionbox}
\textbf{A truck that loses value the way a drug leaves the body.} A delivery truck costs \$32{,}000 new
and loses $18\%$ of its value every year.
\begin{enumerate}[label=(\alph*), leftmargin=1.8em, itemsep=5pt]
  \item Write the model $V(t)$, then find the truck's value after $4$ years. Say which number in your
        model is the $18\%$, and why it is not $0.18$.
  \item A ``half-life'' for a truck would be the time until it is worth half what it cost. Write the
        equation for that, then divide to get it into the form $(0.82)^{t}=\text{something}$.
  \item You know $(0.82)^{3}\approx0.551$ and $(0.82)^{4}\approx0.452$. Between which two whole years
        does the truck lose half its value? Does your answer to (b) have a common base?
  \item According to this model, does the truck ever become worthless? Answer with the range and the
        asymptote --- then say in one sentence why that is a limitation of the \emph{model} rather
        than a fact about trucks.
\end{enumerate}
\ansline{(a) $V(t)=32{,}000(0.82)^{t}$ and $V(4)=\$14{,}467.90$. The $18\%$ is inside the base: losing $18\%$ leaves $82\%$, so $b=1-0.18=0.82$.}
\ansline{\hspace{1.2em} $b=0.18$ would mean the truck keeps only $18\%$ of its value each year --- Lesson 6.1's percent-as-base error.}
\ansline{(b) $32{,}000(0.82)^{t}=16{,}000$, and dividing by $32{,}000$ gives $(0.82)^{t}=0.5$.}
\ansline{(c) Between year $3$ and year $4$, since $0.551>0.5>0.452$. No common base: $0.5$ is not a power of $0.82$, so this joins the running list.}
\ansline{(d) No. The range is $(0,32{,}000]$ with a horizontal asymptote at $V=0$, so the model's value shrinks forever without reaching zero.}
\ansline{\hspace{1.2em} A real truck stops running, gets scrapped, or hits a fixed salvage value --- the exponential is a good description for the early years only.}
\end{extensionbox}

\begin{spiralbox}
\textbf{Coming next --- Lesson 6.5: Exponential Regression \& Choosing a Model.} Today every model was
handed to you: a rate, a half-life, a doubling time. Real data arrives as a pile of points with nobody
to tell you which family it belongs to. Next lesson you bring back Lesson 6.0's fingerprint test ---
constant difference, constant second difference, constant \emph{ratio} --- and use it to decide whether
a data set is linear, quadratic, or exponential before letting technology fit a curve. Then you
interpret the $a$ and the $b$ the calculator hands you, exactly as you did today with $P$ and
$1+\frac{r}{n}$. \textbf{And the list is still standing:} $2^{x}=5$, $132(0.85)^{t}=32$,
$(0.9)^{t}=2.118$, $(1.05)^{t}=2$, and now $(0.82)^{t}=0.5$. Every one has an answer; not one has a
common base. Unit 7 is where you finally write those answers down.
\end{spiralbox}

\begin{teachernote}
\textbf{Problem 1(d) is the one to check first.} The exponent is $365\times4=1460$ and students who
write $4$ or $365$ have lost the sentence the lesson is built on. The balance (\$1144.10) is also only
about a dollar above the same account compounded monthly, which is the Hook restated in homework form
--- ask about it in the warm-up debrief tomorrow.

\textbf{Problems 5(a) and 5(b) are the same error in two different models, and that is the point.}
Tobias keeps an annual base with a monthly count; Marisol keeps a per-half-life base with an hourly
count. Both are ``the base and the exponent stopped agreeing about one period.'' The absurd sizes of
their answers (\$81{,}969 from \$5000; $0.00034$ mg from $90$ mg) are the reasonableness check students
should be running on their own work --- name that habit out loud.

\textbf{Problem 4 rehearses the Section 4 graph with different numbers, and 4(d) is the assessed
sentence.} A student who writes ``no solution'' rather than ``no common base'' has the misconception
that Lesson 7.4 exists to fix; flag those papers and open Lesson 6.5 with thirty seconds on the graph.
The Rule of 72 predicts $12$ years here, which brackets nicely with the graph read --- mention it only
if a student who did Tier E raises it.

\textbf{The Extension is where depreciation and half-life are revealed to be the same object.} Part (c)
adds a fifth equation to the running list, and part (d) is the honest modeling beat: the asymptote is
mathematically correct and physically wrong, which is exactly the kind of judgement Lesson 6.5's
extrapolation discussion needs. Accept any sensible limitation --- salvage value, scrapping, a total
loss --- as long as the range and asymptote are stated first.
\end{teachernote}

\end{document}
